\documentclass[11pt]{article}
\usepackage[margin=1in]{geometry}
\usepackage{amsmath,amssymb}
\usepackage{booktabs}
\usepackage{graphicx}
\usepackage{hyperref}
\usepackage{xcolor}
\usepackage{microtype}
\hypersetup{hidelinks}

\title{TN-Check: Time-Bounded CSL Model Checking of Chemical Master Equations\\ on Tensor-Train State Vectors}
\author{TN-Check authors}
\date{Preprint}

\begin{document}
\maketitle

\begin{abstract}
Time-bounded reachability on the chemical master equation (CME) is a core query
in stochastic systems biology, and probabilistic model checkers answer it by
uniformisation on an explicitly stored state space. The state space of a reaction
network grows as the product of the copy-number ranges of its species, and
explicit and hybrid engines store at least one probability vector over the
relevant states. TN-Check evaluates the time-bounded until operator of CSL
directly on tensor-train (TT) representations of the probability vector and the
CME generator. It applies the absorbing projection of the until operator to the
generator in matrix-product-operator form, runs uniformisation with TT rounding,
and reports the Poisson truncation mass and the finite-state-projection (FSP)
leakage alongside the result. On a shared suite exported with identical FSP
semantics to Storm 1.14 and PRISM 4.10, TN-Check agrees with Storm to within
$3\cdot10^{-9}$ on all cascade models up to $3.4\cdot10^7$ states and on toggle
switches up to $1.6\cdot10^5$ states. On an $n$-stage gene-expression cascade
with 32 copy numbers per species it answers the query for $n=6$ to $n=10$
($1.1\cdot10^{9}$ to $1.1\cdot10^{15}$ states) in 3 to 18 minutes on two CPU
cores and 4~GiB of memory, with every answer inside the 99\% confidence interval
of $10^6$ independent SSA runs. With 16~GiB, Storm does not finish $n=6$ within
an hour and PRISM runs out of memory or time on $n=6$. The TT-CME solver of Ion et
al.\ is faster on the largest cascades and agrees with TN-Check to within
$5\cdot10^{-6}$ there; where an exact reference exists, TN-Check's error is one to
five orders of magnitude smaller than TT-CME's.
\end{abstract}

\section{Introduction}

A stochastic reaction network with $N$ species, each truncated to $d$ copy numbers,
has a finite-state-projection (FSP) state space of $d^N$ states
\cite{munsky2006fsp}. Continuous stochastic logic (CSL) model checking of such a
model asks for probabilities such as $\Pr[\mathit{true}\ \mathrm{U}^{\le T}\ \phi]$,
the probability that a target set is reached by time $T$. Storm
\cite{hensel2022storm} and PRISM \cite{kwiatkowska2011prism} compute this by
uniformisation \cite{baier2003csl,fox1988poisson} on an explicit sparse matrix or
on a hybrid representation that still stores probability vectors explicitly. The
memory cost of these vectors, $8 d^N$ bytes each, bounds the models they can
check.

Tensor trains (TT, also called matrix product states) \cite{oseledets2011tt} store a
$d^N$ vector with $O(N d \chi^2)$ numbers, where $\chi$ is the TT rank. Earlier work
showed that CME solutions of many networks have low TT rank and solved the CME
directly in TT or quantised-TT format
\cite{kazeev2014qtt,dolgov2015state,ion2021ttcme}. The inspected contributions address transient CME solutions and parameter
inference. The candidate distinction here is the CSL until projection and
query interface; TT compression and scalable CME evolution are established
prior work, and a universal absence of earlier TT property checking has not
been established.

\paragraph{Claim.} Time-bounded CSL until on FSP-truncated CMEs can be evaluated in
TT format with accuracy comparable to an explicit model checker on models that
both can handle, and on models whose state space is several orders of magnitude
beyond what explicit and hybrid engines handle in the same memory.

\paragraph{Contributions.}
\begin{itemize}
\item An algorithm for $\Pr[\phi_1\,\mathrm{U}^{\le T}\,\phi_2]$ on TT state vectors:
  the absorbing projection is applied to the generator as a column mask in MPO
  form, uniformisation is accumulated on scalars, and the result carries the Poisson
  tail and the FSP leakage as explicit error terms (Section~\ref{sec:approach}).
\item A PRISM exporter with the same FSP sink semantics as the TT compiler, which
  makes TN-Check, Storm and PRISM answer the identical query on the identical CTMC
  (Section~\ref{sec:impl}).
\item An evaluation against Storm (sparse and hybrid engines), PRISM (explicit,
  sparse, hybrid, MTBDD), the TT-CME solver of Ion et al.\ \cite{ion2021ttcme}, an
  independent scipy reference, and Gillespie simulation \cite{gillespie1977ssa},
  on models from $50$ to $1.1\cdot10^{15}$ states (Section~\ref{sec:eval}).
\end{itemize}

\section{Background and related work}

\paragraph{CSL model checking of CTMCs.} For a CTMC with generator $Q$ (column
convention, $\dot p = Qp$), $\Pr[\phi_1\,\mathrm{U}^{\le T}\,\phi_2]$ equals the
probability mass in $\mathrm{Sat}(\phi_2)$ at time $T$ of the chain in which every
state outside $\mathrm{Sat}(\phi_1)\setminus\mathrm{Sat}(\phi_2)$ is absorbing
\cite{baier2003csl}. Uniformisation writes $e^{Q'T}p_0=\sum_k
\mathrm{Pois}(k;qT)\,P^kp_0$ with $P=I+Q'/q$ and $q$ at least the largest exit rate;
the Poisson series is truncated with Fox--Glynn style bounds \cite{fox1988poisson}.
Storm and PRISM implement this on sparse matrices and on MTBDD-based hybrid
representations. STAMINA \cite{neupane2019stamina} truncates infinite CTMCs on the
fly and bounds the probability of the truncated part; Mikeev et al.\
\cite{mikeev2013dta} check timed-automaton properties on large Markov population
models with on-the-fly state-space truncation. These methods keep an explicit
vector over the retained states.

\paragraph{Tensor methods for the CME.} Kazeev et al.\ \cite{kazeev2014qtt} solved the
CME in quantised TT format; Dolgov and Khoromskij \cite{dolgov2015state} solved it
with a simultaneous space-time TT representation and the AMEn solver
\cite{dolgov2014amen}; Ion et al.\ \cite{ion2021ttcme} use TT for the CME and for
parameter inference and provide the TT-CME code we compare against. Buchholz et
al.\ \cite{buchholz2017compact} study hierarchical-Tucker compression of solution
vectors in Kronecker-structured Markov models. Nicholson and Gingrich
\cite{nicholson2023rare} use tensor networks for rare-event probabilities in
reaction-diffusion CMEs.

\section{Approach}\label{sec:approach}

\paragraph{Generator.} Every reaction $r$ with stoichiometry $v_r$ and a propensity
that factors over species, $a_r(n)=\prod_k f_{r,k}(n_k)$ (mass action, Hill and
Michaelis--Menten kinetics do), contributes a rank-one MPO
$\bigotimes_k S_{r,k}F_{r,k}$, where $F_{r,k}=\mathrm{diag}(f_{r,k})$ and $S_{r,k}$
shifts species $k$ by $v_{r,k}$, plus a rank-one diagonal term $-\bigotimes_k F_{r,k}$.
Shifts that leave the box are dropped from the off-diagonal part but kept on the
diagonal, so probability that leaves the box is lost; this is the FSP sink. The sum
is rounded to a compact MPO.

\paragraph{Absorbing projection.} Let $m=\mathrm{Sat}(\phi_1)\odot(1-\mathrm{Sat}(\phi_2))$.
For axis-aligned atomic propositions both satisfaction sets are rank-one TTs, so
$m$ has rank one and $Q'=Q\,\mathrm{diag}(m)$ has the MPO rank of $Q$. Zeroing columns
makes the states outside $m$ absorbing while keeping the inflow into
$\mathrm{Sat}(\phi_2)$.

\paragraph{Uniformisation in TT format.} With
$q\ge\max_n\sum_r a_r(n)$, computed from the per-species factors as
$\sum_r\prod_k\max f_{r,k}$, and $P=I+Q'/q$, TN-Check iterates
$x_{k+1}=\mathcal{R}_{\chi,\varepsilon}(Px_k)$, where $\mathcal{R}$ is TT-SVD rounding
\cite{oseledets2011tt} to rank at most $\chi$ and relative accuracy $\varepsilon$. Only
scalars are accumulated:
\[
\hat p=\sum_{k=0}^{K} w_k\,\langle \mathrm{Sat}(\phi_2),x_k\rangle,\qquad
\hat m=\sum_{k=0}^{K} w_k\,\langle \mathbf 1,x_k\rangle,\qquad
w_k=\mathrm{Pois}(k;qT),
\]
with $K$ the smallest index whose Poisson right tail is below $10^{-10}$. Inner
products with rank-one masks cost $O(Nd\chi^2)$, so no full vector is formed.

\paragraph{Error terms.} TN-Check reports three quantities with every answer. The
Poisson tail $\tau=1-\sum_k w_k$ bounds the uniformisation truncation, because every
exact iterate has $\|P^kp_0\|_1\le1$. For exact iterates, the deficit $s=1-\hat m-\tau$ bounds the FSP
error from above: the probability of the untruncated CME lies in
$[\hat p,\hat p+s]$ up to the other terms \cite{munsky2006fsp}. The TT rounding error is
estimated by rerunning with $2\chi$ and $\varepsilon/100$ and taking the difference of
the two answers; the checker widens its interval by this amount and returns
INDETERMINATE when the interval contains the threshold. The rounding estimate is
empirical; the Poisson and FSP terms hold for the exact iterates. For an
untruncated nonexplosive CME, the latter requires the same until property
and consistent initial distribution; it concerns paths that exit the box
before termination. Finite-box agreement alone does not validate this
infinite-state approximation. Rounding the
generator and the uniformisation MPO also introduces perturbations. The
reported interval is therefore a numerical diagnostic. A trace consistency
check validates the recorded scalar arithmetic and threshold comparison.

\paragraph{Unbounded until.} For unbounded until TN-Check doubles the horizon $T$ and
uses $\Pr[\mathrm{U}^{\le T}]\le\Pr[\mathrm{U}]\le\Pr[\mathrm{U}^{\le T}]+\langle m,x(T)\rangle+s$,
where $\langle m,x(T)\rangle$ is the mass still in transient states, until the interval
is narrower than a tolerance.

\section{Implementation}\label{sec:impl}

TN-Check is a Python package (numpy, scipy) with a command-line interface that reads
YAML, JSON or SBML networks and CSL formulas. The TT engine is
\path{tn_check/checker/tt_transient.py}. Three components exist only for
validation and share no numerical code with the engine:
an explicit FSP generator assembled state by state from the propensity functions
with scipy \texttt{expm\_multiply} for the until probability
(\path{reference/explicit_cme.py}); a vectorised Gillespie simulator with Wilson and
Clopper--Pearson intervals \cite{wilson1927} (\path{reference/ssa.py}); and a PRISM
exporter (\texttt{export/prism.py}). The exporter writes one bounded integer variable
per species, guards every command by the box, and redirects out-of-box firings to an
absorbing sink in which all counts are reset, so the exported CTMC is the FSP chain
used by the TT compiler; atomic propositions exclude the sink.

\section{Evaluation}\label{sec:eval}

\paragraph{Questions.}
RQ1: Does TN-Check compute the same probabilities as Storm and PRISM?
RQ2: How do time and memory compare with explicit, hybrid and TT baselines?
RQ3: Does TN-Check give correct answers on models that the explicit tools cannot check?

\paragraph{Models and properties.} The suite contains the four example networks of the
repository (birth--death, two-stage gene expression, a toggle switch with Hill
repression, and Michaelis--Menten enzyme kinetics), the toggle switch with boxes of
$200^2$ and $400^2$, and an $n$-stage cascade for $n=2,\dots,10$ in which $S_1$ is
produced at rate 1, $S_i$ is produced with Hill propensity $x_{i-1}^2/(25+x_{i-1}^2)$,
and each species degrades at rate $0.1$, with copy numbers $0,\dots,31$ per species.
Each model has one property $\Pr[\mathit{true}\ \mathrm{U}^{\le T}\ x_k\ge c]$; for the
cascade it is $x_n\ge3$ with $T=10n$.

\paragraph{Setup.} All tools run in separate Modal containers: Storm 1.14.0 (Docker image
\path{movesrwth/storm:1.14.0}) and PRISM 4.10.1 with 1 CPU and 16~GiB (Java heap 12~GiB),
TN-Check and the SSA with 2 CPUs and 4~GiB, TT-CME (commit \texttt{4aebc15}, torchTT 0.5.0,
Crank--Nicolson time stepping, AMEn tolerance $10^{-8}$) with 2 CPUs and 8~GiB.
TT-CME uses a reflecting copy-number boundary, while TN-Check, Storm and PRISM
use an FSP sink; boundary excursions can therefore affect their probability
difference. Time is wall-clock time of the process, memory is its peak resident set. The timeout
is 600~s for the small models and 1200~s for $n=5,6$. The reference value is the scipy
explicit computation where it fits (all models up to $10^6$ states) and Storm sparse
otherwise. TN-Check runs with $(\chi,\varepsilon)=(32,10^{-8})$ and $(64,10^{-10})$.

\paragraph{RQ1: agreement.} Table~\ref{tab:h2h} lists the absolute difference from the
reference. With $\chi=64$, TN-Check is within $7\cdot10^{-11}$ of the reference on
birth--death and the three toggle switches, within $5.4\cdot10^{-9}$ on gene expression,
and within $3\cdot10^{-9}$ on the cascades up to $n=5$. Storm sparse and hybrid agree with
the scipy reference to within $4\cdot10^{-11}$; PRISM agrees to within $6\cdot10^{-8}$ with its
default settings, except on gene expression, where its sparse, hybrid and MTBDD engines
return a value $7\cdot10^{-5}$ below the reference. The enzyme kinetics model has two
conservation laws, and in the species ordering used its distribution needs TT rank 120:
at $\chi=32$ and $\chi=64$ TN-Check is off by $0.25$ and $10^{-2}$, the rank-doubling
estimate reports a difference of $0.24$, and the checker returns INDETERMINATE for
thresholds inside the widened interval; at $\chi=128$ the error is $2\cdot10^{-9}$ and at
$\chi=256$ it is $6\cdot10^{-12}$. Across the 51 models of the repository's model library,
the MPO generator matches propensities evaluated directly at 200 random states of every
model to within $1.6\cdot10^{-13}$, and on the 47 models with at most $1.1\cdot10^6$ states
the TT answer to a reachability property differs from the scipy reference by at most
$4.1\cdot10^{-9}$.

\begin{table}[t]
\centering\small
\resizebox{\textwidth}{!}{\begin{tabular}{lrrrrrrrrr}
\toprule
 & & \multicolumn{2}{c}{Storm sparse} & PRISM sparse & \multicolumn{2}{c}{tn\_check ($\chi{=}64$)} & \multicolumn{2}{c}{TT-CME ($\varepsilon{=}10^{-8}$)} \\
model & box states & s & GB & s & s & abs.\ err & s & abs.\ err \\
\midrule
birth\_death & $5.0\cdot 10^{1}$ & 0.2 & 0.2 & 1.4 & 2.1 & 4e-11 & 10.1 & 3e-6 \\
gene\_expression & $8.0\cdot 10^{3}$ & 0.3 & 0.2 & 2.0 & 9.3 & 5e-9 & 311 & 9e-8 \\
toggle\_switch & $1.0\cdot 10^{4}$ & 0.2 & 0.2 & 1.6 & 2.8 & 2e-11 & 42.3 & 2e-6 \\
toggle\_d200 & $4.0\cdot 10^{4}$ & 0.4 & 0.2 & 2.4 & 42.2 & 4e-11 & 234 & 5e-7 \\
toggle\_d400 & $1.6\cdot 10^{5}$ & 0.8 & 0.2 & 9.2 & 140 & 6e-11 & 1763 & 6e-7 \\
enzyme\_kinetics & $6.4\cdot 10^{5}$ & 0.1 & 0.2 & 1.4 & 26.1 & 1e-2 & 47.6 & 2e-5 \\
cascade\_n2 & $1.0\cdot 10^{3}$ & 0.1 & 0.2 & 1.8 & 2.4 & 5e-11 & 8.6 & 1e-6 \\
cascade\_n3 & $3.3\cdot 10^{4}$ & 0.1 & 0.2 & 1.6 & 2.8 & 1e-9 & 17.3 & 7e-7 \\
cascade\_n4 & $1.0\cdot 10^{6}$ & 1.2 & 0.2 & 12.9 & 22.0 & 2e-9 & 42.2 & 3e-7 \\
cascade\_n5 & $3.4\cdot 10^{7}$ & 45.7 & 2.5 & 349 & 102 & 3e-9 & 66.9 & 2e-7 \\
cascade\_n6 & $1.1\cdot 10^{9}$ & TO & -- & TO & 571 & -- & 153 & -- \\
\bottomrule
\end{tabular}
}
\caption{Wall time (s), peak memory (GB) and absolute error against the reference.
TO: timeout. For $n=6$ the TT error column is empty because no exact reference exists;
see Table~\ref{tab:scale}.}
\label{tab:h2h}
\end{table}

\paragraph{RQ2: cost.} Where the state space fits, Storm sparse is the fastest tool by one
to two orders of magnitude: it builds only the states relevant to the property (for the
cascade, the states with $x_n\le3$) and needs 0.1 to 46~s on all models up to $n=5$.
TN-Check needs 2 to 140~s on the same models. Its cost grows with the number of
uniformisation steps, which is about $qT$ (116 to 10\,665 MPO applications in this suite),
and with $\chi^3$. TT-CME is slower than TN-Check at $\chi=64$ on all models up to $n=4$ and
faster on the cascades from $n=5$ on (Tables~\ref{tab:h2h} and~\ref{tab:scale}): its implicit time
stepping does not depend on $q$, while TN-Check's step count grows with $n$ through both
$q$ and $T=10n$. At AMEn tolerance $10^{-8}$ the TT-CME error against the exact reference is
$9\cdot10^{-8}$ to $1.7\cdot10^{-5}$ across the suite, compared with $2\cdot10^{-11}$ to
$5.4\cdot10^{-9}$ for TN-Check at $\chi=64$ (excluding the enzyme model). These observed accuracy differences are not a matched-accuracy ranking:
the solvers use different boundaries, numerical schemes and stopping criteria.
A matched-boundary, matched-error cost comparison remains open.

\paragraph{RQ3: beyond explicit engines.} For $n=6$ the relevant state space explored by
Storm has $1.3\cdot10^8$ states and $1.2\cdot10^9$ transitions. Storm sparse builds the model
in 568 to 906~s and does not finish the transient analysis within 3\,500~s; its hybrid
engine stops with a full BDD unique table; PRISM explicit runs out of its 12~GiB Java heap,
and PRISM sparse and hybrid do not finish within 1\,200~s. For $n=7$ Storm sparse is still
building the model after 2\,400~s; the relevant state space has about $4\cdot32^6\approx
4.3\cdot10^9$ states, so a single probability vector over it takes 34~GB. TN-Check answers
$n=6,\dots,10$ with $\chi=32$ in 3 to 18 minutes and with $\chi=64$ in 10 to 50 minutes,
using at most 0.7~GB of memory (Table~\ref{tab:scale}). Each answer lies inside the Wilson
99\% interval of $10^6$ SSA runs; the change from $\chi=32$ to $\chi=64$ is at most
$1.9\cdot10^{-5}$; and the $\chi=64$ answers agree with TT-CME at tolerance $10^{-8}$ to within
$5\cdot10^{-6}$.

\begin{table}[t]
\centering\small
\resizebox{\textwidth}{!}{\begin{tabular}{rrlllrrrl}
\toprule
$n$ & box states & tn\_check $\chi{=}32$ & tn\_check $\chi{=}64$ & $|\Delta|$ & s ($\chi{=}32$) & s ($\chi{=}64$) & TT-CME $\varepsilon{=}10^{-8}$ & SSA 99\% CI ($10^6$ runs) \\
\midrule
3 & $3.3\cdot 10^{4}$ & 0.886737 & 0.886737 & 4e-08 & 1.5 & 2.8 & 0.886738 (17.3) & [0.8859, 0.8875] \\
4 & $1.0\cdot 10^{6}$ & 0.809520 & 0.809520 & 1e-07 & 12.3 & 22.0 & 0.809520 (42.2) & [0.8088, 0.8108] \\
5 & $3.4\cdot 10^{7}$ & 0.727022 & 0.727022 & 3e-07 & 59.2 & 102 & 0.727023 (66.9) & [0.7255, 0.7278] \\
6 & $1.1\cdot 10^{9}$ & 0.644761 & 0.644766 & 4e-06 & 178 & 571 & 0.644766 (153) & [0.6434, 0.6459] \\
7 & $3.4\cdot 10^{10}$ & 0.566100 & 0.566107 & 7e-06 & 288 & 860 & 0.566107 (210) & [0.5652, 0.5677] \\
8 & $1.1\cdot 10^{12}$ & 0.492950 & 0.492963 & 1e-05 & 342 & 1366 & 0.492962 (379) & [0.4921, 0.4946] \\
9 & $3.5\cdot 10^{13}$ & 0.426288 & 0.426299 & 1e-05 & 1019 & 2633 & 0.426297 (648) & [0.4243, 0.4268] \\
10 & $1.1\cdot 10^{15}$ & 0.366437 & 0.366456 & 2e-05 & 1077 & 2976 & 0.366451 (790) & [0.3660, 0.3685] \\
\bottomrule
\end{tabular}
}
\caption{Cascade models with 32 copy numbers per species. $|\Delta|$: change between
the $\chi=32$ and $\chi=64$ runs. Times in seconds.}
\label{tab:scale}
\end{table}

\paragraph{Small-model rank stress test.} Six model/horizon cases comprise a
birth--death reachability control, an empty-goal leaking control, a $10\times10$
toggle switch at $T=1,5$, and a $6^3$ three-stage cascade at $T=1,5$. At each
requested rank $1,2,4,8$, the checker reports the solution computed at twice
that rank and widens its interval using the rank difference. An independently
enumerated generator is solved by both sparse and dense matrix exponentials.
The maximum difference between these references is $1.12\cdot10^{-15}$; all 24 reported intervals
contain the finite-box reference within a $10^{-9}$ comparison tolerance.
The largest probability error at requested rank 1 is 0.1116, while at
requested rank 4 or 8 it is below $2.8\cdot10^{-11}$. These finite tests
measure numerical behavior across rank caps; the rounding term remains an
empirical estimate. Raw outputs are in \texttt{results/adversarial\_before.jsonl}.

\paragraph{Finite-box residual witness.} A fresh frozen diagnostic enumerates
an independent substochastic uniformisation matrix $P$ over the box and
records every approximate TT iterate $x_k$. In exact arithmetic, with
$x_0=p_0$, define $d_j=\|x_j-Px_{j-1}\|_1$. Contractivity gives
$\|x_k-P^kp_0\|_1\le\sum_{j=1}^k d_j$ by induction, hence the goal-scalar
error is at most $B=\tau+\sum_{k=0}^K w_k\sum_{j=1}^k d_j$.
This residual includes rounded-operator perturbations because $P$ is
independently enumerated. Twenty runs use boxes of 10, 100 and 216 states,
horizons 1 and 5, and ranks 1, 2, 4, 8. All computed envelopes contain the
sparse-exponential reference within $10^{-12}$; low-rank toggle radii range
from 1.03 to 3.30 and therefore yield no useful probability restriction.
A deliberately injected identity operator gives zero reachability at both
ranks 1 and 2, while the reference is 0.075866; rank agreement misses this
fault and the residual radius is approximately 2. This is an artificial
control, not a failure observed in the unmodified engine. Across all 22 runs,
TT iteration with witness evaluation takes 0.177\,s, witness evaluation
0.029\,s, and explicit transient solves 0.046\,s (one CPU thread; indicative
microbenchmarks). The witness enumerates the box and sacrifices TT memory
scaling. Floating-point residuals and Poisson weights are not outward-rounded,
so this implementation remains a diagnostic rather than a certified enclosure.
Code and every row are in \texttt{benchmarks/scripts/residual\_witness.py} and
\texttt{results/residual\_witness.jsonl}.

\section{Limitations}

The rounding term of the error report is an empirical estimate; a small change
between ranks does not establish a universal error bound. The rounded MPO
perturbation is outside the Poisson and FSP bounds for exact iterates, and the evaluation
covers one scalable family (a feed-forward cascade) beyond the reach of explicit tools;
networks with feedback across many species can have higher TT ranks. Models whose
reachable set is much smaller than the copy-number box, such as networks with
conservation laws, are better served by an explicit engine that explores reachable
states. The uniformisation step count grows with the maximal exit rate over the box,
so stiff networks with large copy-number bounds are expensive. TN-Check handles atomic
propositions on single species and their Boolean combinations; nested probability
operators are not part of this evaluation.

\section{Conclusion}

Applying the until projection to the CME generator in MPO form and running
uniformisation on TT vectors lets a CSL query that explicit engines answer exactly on
$10^7$ states be answered with errors of order $10^{-9}$ on those models and within
Monte Carlo confidence intervals on $10^{9}$ to $10^{15}$ states. All raw outputs, job
definitions and scripts are in the repository.

\bibliographystyle{plain}
\begin{thebibliography}{99}
\bibitem{baier2003csl} C.~Baier, B.~Haverkort, H.~Hermanns, J.-P.~Katoen. Model-checking algorithms for continuous-time Markov chains. \emph{IEEE Trans. Softw. Eng.} 29(6):524--541, 2003. doi:10.1109/TSE.2003.1205180
\bibitem{buchholz2017compact} P.~Buchholz, T.~Dayar, J.~Kriege, M.~C.~Orhan. On compact solution vectors in Kronecker-based Markovian analysis. \emph{Performance Evaluation} 115:132--149, 2017. doi:10.1016/j.peva.2017.08.002
\bibitem{dolgov2014amen} S.~V.~Dolgov, D.~V.~Savostyanov. Alternating minimal energy methods for linear systems in higher dimensions. \emph{SIAM J. Sci. Comput.} 36(5):A2248--A2271, 2014. doi:10.1137/140953289
\bibitem{dolgov2015state} S.~V.~Dolgov, B.~N.~Khoromskij. Simultaneous state-time approximation of the chemical master equation using tensor product formats. \emph{Numer. Linear Algebra Appl.} 22(2):197--219, 2015. doi:10.1002/nla.1942
\bibitem{fox1988poisson} B.~L.~Fox, P.~W.~Glynn. Computing Poisson probabilities. \emph{Commun. ACM} 31(4):440--445, 1988. doi:10.1145/42404.42409
\bibitem{gillespie1977ssa} D.~T.~Gillespie. Exact stochastic simulation of coupled chemical reactions. \emph{J. Phys. Chem.} 81(25):2340--2361, 1977. doi:10.1021/j100540a008
\bibitem{hensel2022storm} C.~Hensel, S.~Junges, J.-P.~Katoen, T.~Quatmann, M.~Volk. The probabilistic model checker Storm. \emph{Int. J. Softw. Tools Technol. Transf.} 24:589--610, 2022. doi:10.1007/s10009-021-00633-z
\bibitem{ion2021ttcme} I.~G.~Ion, C.~Wildner, D.~Loukrezis, H.~Koeppl, H.~De Gersem. Tensor-train approximation of the chemical master equation and its application for parameter inference. \emph{J. Chem. Phys.} 155:034102, 2021. doi:10.1063/5.0045521
\bibitem{kazeev2014qtt} V.~Kazeev, M.~Khammash, M.~Nip, C.~Schwab. Direct solution of the chemical master equation using quantized tensor trains. \emph{PLoS Comput. Biol.} 10(3):e1003359, 2014. doi:10.1371/journal.pcbi.1003359
\bibitem{kwiatkowska2011prism} M.~Kwiatkowska, G.~Norman, D.~Parker. PRISM 4.0: Verification of probabilistic real-time systems. In \emph{CAV 2011}, LNCS 6806, pp.~585--591, 2011. doi:10.1007/978-3-642-22110-1\_47
\bibitem{mikeev2013dta} L.~Mikeev, M.~R.~Neuh\"au{\ss}er, D.~Spieler, V.~Wolf. On-the-fly verification and optimization of DTA-properties for large Markov chains. \emph{Form. Methods Syst. Des.} 43:313--337, 2013. doi:10.1007/s10703-012-0165-1
\bibitem{munsky2006fsp} B.~Munsky, M.~Khammash. The finite state projection algorithm for the solution of the chemical master equation. \emph{J. Chem. Phys.} 124:044104, 2006. doi:10.1063/1.2145882
\bibitem{neupane2019stamina} T.~Neupane, C.~J.~Myers, C.~Madsen, H.~Zheng, Z.~Zhang. STAMINA: STochastic approximate model-checker for INfinite-state analysis. In \emph{CAV 2019}, LNCS 11561, pp.~540--549, 2019. doi:10.1007/978-3-030-25540-4\_31
\bibitem{nicholson2023rare} S.~B.~Nicholson, T.~R.~Gingrich. Quantifying rare events in stochastic reaction-diffusion dynamics using tensor networks. \emph{Phys. Rev. X} 13:041006, 2023. doi:10.1103/PhysRevX.13.041006
\bibitem{oseledets2011tt} I.~V.~Oseledets. Tensor-train decomposition. \emph{SIAM J. Sci. Comput.} 33(5):2295--2317, 2011. doi:10.1137/090752286
\bibitem{wilson1927} E.~B.~Wilson. Probable inference, the law of succession, and statistical inference. \emph{J. Amer. Statist. Assoc.} 22(158):209--212, 1927. doi:10.1080/01621459.1927.10502953
\end{thebibliography}

\end{document}
